\chapter{設計方針}

本章では本研究の設計方針を述べる．3.1 節で全体の構成を示し，3.2 節で確信度の引き出し方式の設計，3.3 節でデータセットの選定，3.4 節で評価指標の定義を述べる．3.5 節で実験条件を，3.6 節で設計上の判断とその根拠を整理し，3.7 節で本章をまとめる．

\section{全体構成}

本研究の評価は，図 3.1 に示す流れで行う．

\begin{figure}[htbp]
\centering
\begin{tabular}{c}
\fbox{\parbox{0.75\linewidth}{\centering 問題データ（問題文・正解・回答形式）}} \\
$\downarrow$ \\
\fbox{\parbox{0.75\linewidth}{\centering プロンプト生成 …… 引き出し方式 × 問題形式でテンプレートを選択}} \\
$\downarrow$ \\
\fbox{\parbox{0.75\linewidth}{\centering モデルへの問い合わせ …… 温度 0 で API を呼び出す}} \\
$\downarrow$ \\
\fbox{\parbox{0.75\linewidth}{\centering 応答の解析 …… 回答と確信度を抽出し，確信度を 0〜1 に正規化}} \\
$\downarrow$ \\
\fbox{\parbox{0.75\linewidth}{\centering 正誤判定 …… 回答形式に応じて正解と照合}} \\
$\downarrow$ \\
\fbox{\parbox{0.75\linewidth}{\centering 指標の計算 …… ECE，Brier Score，AUROC，Murphy 分解}} \\
$\downarrow$ \\
\fbox{\parbox{0.75\linewidth}{\centering 統計的検証 …… 問題の再抽出による信頼区間}} \\
\end{tabular}
\caption{評価の処理フロー}
\label{fig:3.1}
\end{figure}

この流れのうち，確信度の引き出し方式によって変わるのはプロンプト生成と応答の解析の 2 箇所のみである．それ以降の処理はすべての方式で共通とする．方式間の差が処理の違いに由来しないようにするためである．

\section{確信度の引き出し方式}

本研究では 3 つの方式を比較する．いずれもモデルの重みには触れず，プロンプトのみで確信度を引き出す．

\subsection{Verb.1S（回答と確信度の同時表明）}

回答と確信度を 1 回のやりとりで同時に答えさせる方式である．Xiong ら \cite{xiong2024llmuncertainty} が基本形として扱っている方式に相当する．プロンプトの構成は以下の通りである．

\begin{quote}
\begin{verbatim}
以下の数学の問題を読み，回答と，その回答が正しいと思う確率（確信度）を
答えてください．確信度は 0.0（まったく自信がない）から 1.0（完全に確信して
いる）の数値で表してください．

問題: {問題文}

以下の形式で回答してください．形式以外の出力はしないでください．
回答: [数値のみ]
確信度: [0.0〜1.0の数値]
\end{verbatim}
\end{quote}

この方式の利点は，API の呼び出しが 1 回で済むため費用と時間が小さいことである．一方で，回答を生成する過程と確信度を見積もる過程が同一の生成の中で混在するため，確信度が回答の内容を踏まえたものになっていない可能性がある．Seo ら \cite{seo2025advice} が指摘する回答非依存性の問題が生じうる．

\subsection{Verb.2S（2 段階での確信度表明）}

回答を得たあと，そのやりとりを保持したまま別のターンで確信度を尋ねる方式である．Tian ら \cite{tian2023justask} の 2 段階方式に相当する．第 1 ターンでは回答のみを求め，第 2 ターンで以下のように尋ねる．

\begin{quote}
\begin{verbatim}
あなたは先ほどの問題に対して「{第1ターンの回答}」と回答しました．
この回答が正しい確率を 0.0（まったく自信がない）から 1.0（完全に確信して
いる）の数値で答えてください．

以下の形式で回答してください．形式以外の出力はしないでください．
確信度: [0.0〜1.0の数値]
\end{verbatim}
\end{quote}

第 2 ターンのプロンプトに自分の回答を明示的に含めることで，確信度の推定が回答内容に依存するよう促すことを狙う．これは 2.2.4 項で述べた ADVICE の診断的知見を動機とした設計であり，狙いどおりの効果が得られるかは検証の対象である．なお，API の呼び出しが 1 問あたり 2 回になるため費用と所要時間は増える．実際の増加分はトークン数や応答長に依存するため，本研究では実測していない．

\subsection{Ling.1S（言語表現による確信度表明）}

確信度を数値ではなく，あらかじめ定めた 5 段階の言語表現から選ばせる方式である．Tian ら \cite{tian2023justask} の言語表現方式に相当する．各表現は表 3.1 に示す数値に対応づけて集計する．

\begin{table}[htbp]
\centering
\small
\caption{言語表現と対応する数値}
\label{tab:3.1}
\begin{tabular}{|l|l|}
\hline
言語表現 & 数値 \\
\hline
ほぼ確実 & 0.95 \\
かなり自信がある & 0.80 \\
どちらともいえない & 0.50 \\
あまり自信がない & 0.25 \\
ほとんどわからない & 0.05 \\
\hline
\end{tabular}
\end{table}

対応づける数値は，5 段階を 0 から 1 の範囲におおむね等間隔で配置しつつ，両端を 0 と 1 にしない値とした．言語表現は程度を表すものであり，「ほぼ確実」を確率 1，すなわち誤る可能性がないという表明として扱うのは強すぎると判断したためである．この割り当ては暫定的なものであり，値を変えた場合に指標がどの程度動くかは別途確認する必要がある（3.6 節）．

この方式は利用者への提示という観点で意味がある．一方，段階数が 5 に限られるため確信度の解像度が落ちる．また数値への対応づけを研究者が決めている以上，その割り当て自体が結果に影響する．この点は本研究の限界として第 6 章で述べる．

\subsection{方式の対応関係}

3 方式は，表 3.2 に示す 2 つの軸の組み合わせとして整理できる．

\begin{table}[htbp]
\centering
\small
\caption{3 方式の設計軸}
\label{tab:3.2}
\begin{tabular}{|l|l|l|}
\hline
方式 & やりとりの回数 & 確信度の形式 \\
\hline
Verb.1S & 1 回 & 数値（連続） \\
Verb.2S & 2 回 & 数値（連続） \\
Ling.1S & 1 回 & 言語表現（5 段階） \\
\hline
\end{tabular}
\end{table}

ただし，これは 2 つの要因を厳密に統制した実験ではない．数値と言語表現では段階数や値への対応づけも同時に変わり，1 回と 2 回のやりとりでは回答を生成する際の指示や文脈も変わる．また 2 回のやりとりと言語表現を組み合わせた条件を設けていないため，2 つの要因の交互作用は評価できない．本研究が比較するのは，実際に使われうる 3 つのプロンプト条件の総合的な差である．差の原因を要因ごとに切り分けるには，段階数をそろえた数値条件を加えるなどの設計が要る．

\section{データセット}

\subsection{選定の要件}

較正の評価に用いるデータセットには 3 つの要件がある．

第一に，正誤が機械的に判定できることである．較正の評価には各問題の正誤が必要であり，判定に人手や別のモデルを介すると，判定の誤りが較正の推定に混入する．

第二に，問題数が十分にあることである．較正指標は確信度を区間に分割して計算するため，各区間に一定数の標本が必要になる．

第三に，日本語として自然であることである．2.3 節で述べたとおり，MlingConf は英語の問題を翻訳したものを用いている．日本語タスクでの較正を論じるなら，日本語として自然な問題が望ましい．

\subsection{予備実験で用いるデータセット}

予備実験では日本語版 MGSM \cite{shi2023mgsm} を用いた．MGSM は算数の文章題からなるGSM8K \cite{cobbe2021gsm8k} の一部を人手で多言語に翻訳したもので，日本語版は 250 問である．

MGSM を選んだ理由は第一・第二の要件を満たすためである．解答が数値であるため正誤判定が厳密に行え，判定の曖昧さが入り込まない．また，翻訳が機械ではなく人手によるため，訳文の品質が一定している．

一方で第三の要件は満たしていない．MGSM は英語で作成された問題を翻訳したものである．日本語で入力し日本語で答えさせている以上，日本語入力に対する較正の測定にはなっているが，対象は算数の文章題に限られ，日本固有の知識や文脈を含む課題には及ばない．したがって予備実験の結果を，日本語のタスク一般に広げて論じることはできない．この制約は第 6 章で述べ，今後は別のデータセットによる追試を行う．

\subsection{今後用いるデータセット}

日本語で作成されたデータセットとして，JGLUE \cite{kurihara2022jglue} に含まれるJCommonsenseQA と，JMMLU を検討している．前者は日本語で作成された常識推論の多肢選択式データセットである．後者は幅広い分野の知識を問う多肢選択式であり，英語からの翻訳による科目と日本固有の内容を扱う科目の双方を含むため，使用する範囲を明示して用いる必要がある．

多肢選択式は正誤判定が機械的に行える一方，選択肢が与えられるため無作為に答えても一定の正答率が得られる．この下限が確信度の分布に影響する可能性があり，自由回答式の MGSM と併せて比較する必要がある．また JCommonsenseQA は選択肢が 5 つであり，現在の実装が想定する 4 つの記号では足りない．実験前に選択肢数に応じて記号を生成し，記号と本文と正解の対応を確認する必要がある．

\section{評価指標}

以下，問題数を $N$，$i$ 番目の問題に対するモデルの確信度を $c_i \in [0,1]$，回答が正解なら $y_i = 1$，不正解なら $y_i = 0$ とする．

\subsection{ECE}

確信度の値域 $[0,1]$ を $M$ 個の区間 $B_1, \dots, B_M$ に等分し，各区間に含まれる問題の集合を同じ記号で表す．ECE は次式で定義される．

\begin{equation}
\mathrm{ECE} = \sum_{m=1}^{M} \frac{|B_m|}{N} \bigl| \mathrm{acc}(B_m) - \mathrm{conf}(B_m) \bigr|
\end{equation}

ここで $|B_m|$ は区間 $B_m$ に含まれる問題数，$\mathrm{acc}(B_m)$ は区間内の正答率，$\mathrm{conf}(B_m)$ は区間内の平均確信度である．すなわち

\begin{equation}
\mathrm{acc}(B_m) = \frac{1}{|B_m|}\sum_{i \in B_m} y_i, \qquad
\mathrm{conf}(B_m) = \frac{1}{|B_m|}\sum_{i \in B_m} c_i
\end{equation}

である．ECE は 0 以上 1 以下の値をとり，小さいほど較正が良い．本研究では$M = 10$ とする．

\subsection{Brier Score}

Brier Score は確信度と正誤の二乗誤差の平均である．

\begin{equation}
\mathrm{BS} = \frac{1}{N} \sum_{i=1}^{N} (c_i - y_i)^2
\end{equation}

0 以上 1 以下の値をとり，小さいほど良い．区間分割を必要としない点で ECE と異なる．

\subsection{Murphy 分解}

Brier Score は次の 3 成分に分解できる \cite{murphy1973partition}．

\begin{equation}
\mathrm{BS} = \mathrm{REL} - \mathrm{RES} + \mathrm{UNC}
\end{equation}

確信度がとりうる値を $K$ 個の区間に分け，区間 $k$ に含まれる問題数を $n_k$，区間内の平均確信度を $\bar{c}_k$，区間内の正答率を $\bar{y}_k$，全体の正答率を $\bar{y}$ とすると，各成分は次式で与えられる．

\begin{equation}
\mathrm{REL} = \frac{1}{N} \sum_{k=1}^{K} n_k (\bar{c}_k - \bar{y}_k)^2
\end{equation}

\begin{equation}
\mathrm{RES} = \frac{1}{N} \sum_{k=1}^{K} n_k (\bar{y}_k - \bar{y})^2
\end{equation}

\begin{equation}
\mathrm{UNC} = \bar{y}(1 - \bar{y})
\end{equation}

REL（信頼度）は，各区間における確信度と正答率のずれを表し，小さいほど良い．RES（識別度）は，各区間の正答率が全体の正答率からどれだけ離れているかを表す．正答しやすい問題と誤答しやすい問題を確信度で区別できているほど大きくなり，大きいほど良い．UNC（不確実度）は，評価対象となる正誤ラベルの平均$\bar{y}$ のみで決まる．

この分解の扱いについて，3 点を明示しておく．

第一に，上式の右辺は，各問題の確信度をその区間の平均 $\bar{c}_k$ に置き換えた予測に対する Brier Score に対応する．個々の確信度を用いた3.4.2 項の $\mathrm{BS}$ とは一般に一致せず，等号が成り立つのは確信度が区間ごとに 1 つの値しかとらない場合である．本論文で $\mathrm{BS}$ として報告する値は 3.4.2 項の定義によるものであり，分解の右辺との差は第 5 章で併記する．

第二に，REL は各区間の正答率 $\bar{y}_k$ を含む．したがって UNC を切り離すことは，REL を正答率から独立にすることを意味しない．区間が 1 つしかない場合，REL は $(\bar{c} - \bar{y})^2$ となり，平均確信度と平均正答率の差そのものになる．

第三に，UNC は $\bar{y}$ で決まるが，$\bar{y}$ は各方式の回答から定まるため方式ごとに異なりうる．また $p$ と $1-p$ は同じ $p(1-p)$ を与えるので，正答率が異なれば UNC も必ず異なるとは限らない．

以上より，本研究は Murphy 分解を，較正と正答率を切り分ける手段としてではなく，Brier Score の内訳を把握するための補助的な情報として用いる．較正の違いのみを取り出すには，回答を固定したうえで確信度の引き出し方だけを変える比較が必要であり，これは本論文の範囲外である．

\subsection{AUROC}

AUROC は確信度を判別のスコアとみなし，正答を陽性，誤答を陰性としたときのROC 曲線の下側面積である．確信度の絶対的な水準ではなく順序のみを評価する．0.5 が無作為な判別に相当する．

\subsection{指標を併用する理由}

4 つの指標はそれぞれ異なる側面を測る．ECE は較正の水準を，AUROC は識別の順序を，Brier Score は両者を合わせた総合的な予測の質を測る．Murphy 分解は区間に集約した予測について，その内訳の目安を与える．

単一の指標では結論を誤る場合がある．たとえば，すべての問題に同じ確信度を答えるモデルは，その値を全体の正答率に一致させれば ECE を 0 にできるが，どの問題が誤りかは全く示せない．このとき AUROC は 0.5 となり，RES は 0 となる．逆に，確信度の順序は正しいが値が一律に高すぎるモデルは，AUROC が高い一方で ECE が悪くなる．複数の指標を併記することで，このような状況を区別できる．

\section{実験条件}

予備実験の条件を表 3.3 に示す．

\begin{table}[htbp]
\centering
\small
\caption{予備実験の条件}
\label{tab:3.3}
\begin{tabular}{|l|l|}
\hline
項目 & 設定 \\
\hline
対象モデル & claude-sonnet-4-6 \\
温度 & 0.0 \\
データセット & 日本語版 MGSM（250 問） \\
問題形式 & 自由回答（数値） \\
引き出し方式 & Verb.1S，Verb.2S，Ling.1S \\
試行回数 & 各方式 1 回 \\
評価指標 & ECE（10 区間），Brier Score，AUROC，Murphy 分解 \\
\hline
\end{tabular}
\end{table}

温度を 0 とするのは，同一の問題に対する出力の揺らぎを抑えるためである．出力が確率的であっても較正そのものは測定できるが，その場合は反復実験によって変動を評価する必要がある．予備実験では条件を絞るため温度 0 を用いた．なお温度 0 は出力が完全に同一になることを保証するものではない．

対象モデルを 1 つに絞ったのは，予備実験の目的が評価基盤の妥当性の確認と方式間の差の把握にあるためである．モデル間の比較は今後の課題とする．

\section{設計上の判断とその根拠}

本章で述べた設計のうち，判断を要した点を表 3.4 に整理する．

\begin{table}[htbp]
\centering
\small
\caption{設計上の判断}
\label{tab:3.4}
\begin{tabular}{|l|l|l|}
\hline
判断 & 選んだ方針 & 根拠 \\
\hline
確信度の取得法 & プロンプトのみ（重みを更新しない） & 商用 API のモデルに適用でき，追加の学習費用が不要（2.4.3 項） \\
比較する方式数 & 3 方式 & 実際に使われうる条件を優先．要因の統制は行っていない（3.2.4 項） \\
主要な評価指標 & ECE・Brier・AUROC を併記し，Murphy 分解を補助的に示す & 単一の指標では較正と識別を区別できない．分解も較正と正答率を切り分けはしない（3.4.3 項，3.4.5 項） \\
予備実験のデータ & MGSM（翻訳データ） & 正誤判定の厳密さを優先．対象が算数の文章題に限られる点は限界とする（3.3.2 項） \\
言語表現の数値化 & 0.05〜0.95 の 5 段階 & 「ほぼ確実」を確率 1 として扱わないための暫定値．感度は要確認（3.2.3 項） \\
温度 & 0.0 & 出力の揺らぎと方式の効果を混同しないため（3.5 節） \\
\hline
\end{tabular}
\end{table}

\section{本章のまとめ}

本章では，3 つの確信度引き出し方式，データセットの選定基準，4 つの評価指標の定義を述べた．評価指標については，ECE の値が確信度の分布に左右されること，および Murphy 分解が較正と正答率を切り分ける手段ではないことを明示し，複数の指標を併記する方針をとった．また，3 方式の比較が要因を統制した実験ではないこと，予備実験で用いる MGSM の対象が算数の文章題に限られることを述べた．次章では，この設計に基づく評価基盤の実装について述べる．
